\section{Conclusion}
\label{sec:conclusion}

A serialization format does not have to choose between being a row format
that is cheap for a small message and a columnar one that is compact for a
table. \colbin{} is both: it decides per slice, at encode time, with a
threshold of eight rows, and applies to the columns a small set of integer
transforms and byte-aligned bit-packing that cost little on a short column
and pay increasingly on a long one. It needs no schema file and no generated
code, and a message that embeds its schema can be read as JSON by a reader
that has never seen its type.
The evaluation shows where this pays and where it does not, on data the
author did not design, and the artifact lets a reader check each number.

\section{Code and data availability}
\label{sec:availability}

The benchmarks, dataset scripts, checksums, raw results and the source of
this paper are available at the tag \code{paper-v1} of its repository:
\begin{itemize}
\item repository: \url{https://github.com/ivanjoz/colbin-benchmarks/tree/paper-v1}
\item commit of the tag: \code{625207318491c32bf36d2894117e6a22697443ba}
\item archive: \url{https://doi.org/10.5281/zenodo.24261124}
\item main machine measured at: \code{\result{machine.commit}},
  \colbin{} \result{machine.colbin-version}
\item second machine measured at: \code{\result{laptop.machine.commit}},
  \colbin{} \result{laptop.machine.colbin-version}
\item \colbin{}: \url{https://github.com/ivanjoz/colbin}
\end{itemize}
The raw results of both machines are in the repository,
\code{results/paper-v1} and \code{results/paper-v1-laptop}.
Running \code{./reproduce.sh} at the tag downloads and verifies the
datasets, runs every measurement, and regenerates every table and figure of
this paper; the tables and the numbers quoted in the text are generated
files, not transcriptions.
